\subsection{Reusable classical readout with local feedback}
\label{sec:source-feedback-transport}

Destructive path reconstruction can be replaced by reusable transport when
local feedback is admitted explicitly. Each participating patch has a
mutable scratch port, a protected zero record and immutable scalar version
records. In addition to the canonical pair mean, allow a patch to export
one of its own records into its own scratch port, reset that port from its
own zero record, and archive twice its own port readout. These local
copy/reset and decoding operations are supplied classical software semantics;
the pair-mean law alone does not select them.

For a support seam $u\!\sim\!v$ and a version value $x$ held at $u$,
perform the following six events: export $x$ into $u$, reset $v$ to zero,
average the pair, archive twice the readout at $v$, and reset $u$ and $v$
from their respective local zeros. The final archive at $v$ equals $x$,
both scratch ports return to zero, and no protected version is changed.
Every local operation reads and writes only its own patch; the mean is the
sole operation spanning the seam.
Counts include identity mean evaluations for a zero payload; they assert
no strict repair progress. All means in the retained episodes are active.

\begin{theorem}[Reusable version transport and semantic order]
\label{thm:source-feedback-transport}
On a supplied connected support, the stated local operations transport an
immutable scalar version along any finite path of $d$ seams using $d$ means,
$d$ exports, $d$ captures and $3d$ resets. The transport uses exactly $6d$
events, $7d$ scalar reads and $7d$ scalar writes, with non-interleaved
six-event hops. It remains valid under
repeated reads, reuse of scratch ports and arbitrary changes of logical
values at their publication, with fixed paths, schedule and read interfaces.

For a finite acyclic scalar read program with total local commits,
compile each requested read to
such a path and let every logical commit consume only its delivered records
and its stated local inputs. The induced semantic read-from order on logical
commits is exactly the program's read-from order. All intermediate events
remain in the refinement.
\end{theorem}
\begin{proof}
The received value is $2(x+0)/2=x$, independently of the previous scratch
values. Induction along the path transports the same immutable version.
Count the six events and their explicit interfaces: the mean has two reads
and two writes; each of the other five events has one of each.

Label each payload event by its originating logical writer and label local
initialization and constant resets by a bottom symbol. A semantic edge out
of a constant has no logical ancestor; other transport edges retain their
source label until a logical commit consumes the payload. Every such edge
therefore maps to equality or a logical read path. Conversely, every
requested read has its retained export--mean--capture path from its writer
to its consumer. Path composition gives both inclusions. Reset writers
depend only on protected zeros, so a prior scratch payload creates no
additional semantic ancestor. The argument holds for every choice of scalar
values and hence commutes with the stated interventions.
\end{proof}

This is a projection of semantic data dependencies. It does not identify
the projected order with physical causal order: shared-resource execution,
scheduler/control order, write-after-write constraints and custody-chain
sequencing remain separate.
Version identity is bound by the fixed compiled schedule and authenticated
writer chain; a physical address/header channel is not supplied.
Feedback can change the live loads used by a carrier's instantaneous
readback; maintained addresses are separate protected records here.

\begin{theorem}[Additive error under restored transport]
\label{thm:source-feedback-error}
For each hop, bound the export, receiver reset, receiver mean coordinate,
readback and decoder errors by $E,Z,M,R,D\geq0$, respectively. Include
archive retention/read/write error in $E$. If the initial archive error is
at most $E_0$, the final depth-$d$ payload error satisfies
\[
 |\widehat x-x|\leq E_0+d(E+Z+2M+2R+D).
\]
The bound is sharp when the stated additive errors may be chosen separately.
\end{theorem}
\begin{proof}
Writing one hop's errors as $e,z,m,r,a$, the receiver obtains
\[
 2\left(\frac{\widehat x+e+z}{2}+m+r\right)+a
 =\widehat x+e+z+2m+2r+a.
\]
The incoming error has coefficient one. Sum over the hops and apply the
triangle inequality. Equal-sign extremal errors attain the bound.
Fresh export and receiver reset precede every mean, so cleanup errors from
an earlier hop are overwritten before they can enter a later payload.
\end{proof}
The noise limits are hypotheses, not measurements of hardware. The exact
event receipt introduces no noisy dynamics or physical precision guarantee.

The independently replayed W12 construction in
\texttt{code/source\_feedback\_transport/} retains 16 complete episodes,
740 seam hops and 4,928 events at path depths $1,4,12,24$. Its routes are
reused and intersect; an intervened branch leaves subsequent reads of an
older version unchanged. At depth 24, five requests use 108 hops and 706
total events, including initialization and logical commits, with 141
protected scalar registers and 25 scratch registers. Version identifiers,
event metadata and arithmetic workspace are additional storage costs.
The Lean module \texttt{SourceFeedbackTransport} proves exact hop and route
identities, error bounds and the finite semantic-projection criterion; the
verifier instantiates that criterion from every consumed writer.

This supplies a classical implementation of declared remote scalar reads.
The complete scalar field execution is not compiled into these episodes.
The intended field update, source selection of feedback and read rules,
quantum operations, outcome maps and physical clock remain additional inputs.
A separate compiler census at $q=13,21$ consumes the earlier declared-family
read counts. It does not execute the complete routed histories at those
sizes, and its unicast cost calculation excludes alternative multicast
compilers.
